\section{Implementation sequence and default admission}
\label{sec:plan}
\subsection{Recommended order}
The four requested areas are substantial formulation and validation
programs. They are not four independent switches at the end of the
existing BVP. The recommended order is to define shared contracts,
complete a smooth forward/derivative chain, add consistent frozen-gas
thermodynamics and engine performance coupling, then expand physical
regimes with explicit applicability.

A bounded attached calorically perfect default can be assessed before
all reacting/separated capabilities exist, provided its user-facing
scope and refusal behavior are explicit. A default intended for realistic
hot-gas design may require the frozen-property stage before admission.
No universal default is justified by the present single-context result.

\begingroup\small
\begin{longtable}{@{}P{0.14\linewidth}P{0.39\linewidth}P{0.39\linewidth}@{}}
\toprule\textbf{Milestone / depends on} & \textbf{Bounded implementation task} &
\textbf{Acceptance evidence}\\\midrule
\endhead
C0 / current code &
Specify gas/geometry/flow domain and one flow/performance result.
Extend typed invalid, incomplete, unsupported-regime and gradient-unavailable
outcomes through engine/NLP/export. &
Contract roundtrips; invalid provider cannot influence an accepted
performance result; no unreported TOP substitution; units and source identity.\\
N0 / C0 &
Make general-wall host unit processes and topology an explicit primal
reference. Bind wall, axis, parents, starting line and exit fluxes to
the same state. &
Matched NASA controls, simple exact/constructed flows, grid convergence,
physical strip and independent momentum/mass/energy checks.\\
N1 / N0 &
Complete the differentiable forward-wall operator, including start data,
intersections, reflections, field projection and total boundary dependence. &
Converged primal parity; directional Taylor tests through supported
design variables; topology-event and conditioning diagnostics.\\
T0 / C0 &
Expose the existing frozen thermodynamics through a consistent property
and derivative contract; define transport separately. &
Constant-\(c_p\) limit; EOS/enthalpy/entropy identities; property-grid
refinement; independent reference states and explicit domain refusal.\\
T1 / N0, T0 &
Derive and implement smooth fixed-composition variable-property
characteristics and gas-consistent sonic/transonic starting data. &
Calorically perfect reduction; Q1D and smooth-flow checks; independent
Euler field/performance comparisons; no local-\(\gamma\) substitution shortcut.\\
M0 / C0, accepted N/T branch &
Feed accepted aerodynamic performance into throat/mass-flow residuals,
then local wall state into thermal/geometry/feed/mass consumers. &
Intentional perturbations affect the correct outputs and constraints;
full-engine gradient tests; consistent \(c^*\)/loss ledger and model identity.\\
V0 / accepted core, T0, thermal interface &
Select and implement attached compressible nozzle BL closures and
core-displacement/thermal coupling within a declared thin-layer domain. &
Pressure, heat, drag/displacement comparisons on source-matched cases;
transition sensitivity; converged interface; thick/separated-domain refusal.\\
S0 / T0 &
Implement conservative shock jumps with explicit gas/front model,
normal orientation, entropy admission and compressive branch selection. &
Analytic perfect-gas limits; common-EOS general-gas reference states;
signed jump/energy residuals and rejected trivial root.\\
S1 / N0, S0 &
Add rotational downstream state, fitted-front degrees of freedom and
selected shock/wall/axis topologies. &
A matched shock-location/field reference, upstream/downstream
conservation and entropy transport; explicit mixed-type limitation.\\
V1 / V0, S1, selected mixed-type method &
Connect resolved wall states to scoped onset screening and introduce
separated-region/reattachment branch treatment where required. &
Geometry/thermal/history-matched FSS/RSS data; separation definition,
pressure recovery and uncertainty; no inferred side-force time history.\\
T2 / T0, equilibrium backend evidence &
Add a consistent equilibrium EOS, composition/acoustic derivatives and
its accepted smooth or shock branch. &
Element/mass/energy closure, independent equilibrium states, acoustic
derivative checks, reference-front limits and derivative-domain tests.\\
T3 / T0, N0, corroborated reacting derivation &
Add finite-rate species and selected mode histories, stiff source
integration and consistent characteristic/sonic initialization. &
Resolve Young source discrepancy; chemistry conservation; homogeneous
reactor/relaxation and nozzle cases; time/space/kinetics sensitivity.\\
A0 / each intended admitted branch &
Establish accuracy, boundary identity, cost, gradient validity and
parameter-domain evidence; test full-engine reoptimization/refinement. &
Independent holdouts and regime edges; restart robustness; runtime/memory
profiles; candidate replay and explicit unsupported states.\\
D0 / C0, N1, M0, A0 and relevant physics &
Promote a named supported branch to default, preserving explicit model
selection, compatibility and qualification status. &
Documented domain and release criteria; same-context CLI/API/MDO
replay; no claim that unimplemented regimes were admitted.\\
\bottomrule
\end{longtable}
\endgroup

The dependency table deliberately permits thermodynamic, attached-layer
and shock-jump research to proceed in parallel. Finite-rate implementation
need not wait for every equilibrium software feature, but it does require
its own corroborated equations and appropriate equilibrium/frozen
reference comparisons. Separated/asymmetric physics is not a dependency
for admitting a strictly attached branch; it is a dependency for making
corresponding off-design claims.

\subsection{Module boundaries and interfaces}
Proposed boundaries are a thermodynamic state/derivative provider; a
flow topology/region representation; continuous MOC unit processes;
shock-jump/front operators; an attached BL/thermal interface; a
separated/mixed-type analysis adapter; and an engine aerodynamic
provider. These are architectural responsibilities, not newly created
production modules.

Keep source-derived equations separate from discretization, nonlinear
solver settings, acceptance policy and fidelity dispatch. Preserve the
existing interfaces through compatible adapters where feasible.
Changing the independent state from \(M,\theta\) to velocity,
thermodynamic state and transported invariants may justify a new versioned
model contract rather than silently changing the meaning of old fields.

The existing Rao design solver can remain a source-controlled reference
and inverse-design option while the forward-wall operator serves engine
MDO. Variable-gas/viscous optimum-design conditions require their own
derivation; solving a more complete forward field does not automatically
generalize the original Rao variational theorem.

\subsection{Release gates and numerical policies}
For each proposed default domain, establish:
\begin{enumerate}
\item governing equations, source conventions, supported topology and
boundary conditions;
\item data identity and EOS/transport/chemistry validity, including
composition and station conventions;
\item actual completed point, global, interface and outer engine solves;
\item conservation, characteristic compatibility and physical mesh checks;
\item field/output convergence under appropriate spatial, property and
nonlinear refinement;
\item independent model-appropriate references and experiments with
traceable reductions and uncertainty;
\item total-gradient validity, conditioning and event behavior;
\item practical runtime/memory and documented initialization/failure behavior;
\item same-model engine reoptimization and reproducible final adjudication.
\end{enumerate}
Acceptance numbers must come from an explicit engineering accuracy and
uncertainty requirement. The historical MOC refinement limits are not
universal tolerances for shocks, heat flux, separation or chemistry.
New thresholds should be set before inspecting a proposed implementation's
failure distribution.

No calendar estimate is assigned from source reading alone. The largest
uncertainties are closure/data availability, transonic initialization,
stable nonlinear coupling, differentiability at events and the actual
cost of accepted multi-physics solves.

\section{Validation assets already in the repository}
\label{sec:validation}
An available PDF is a candidate evidence source, not a ready-to-run
benchmark. Reconstruct geometry, gas/composition, wall thermal condition,
pressure/history, instrumentation, data reductions and uncertainty before
assigning a numerical pass/fail threshold.

\begin{longtable}{@{}P{0.23\linewidth}P{0.34\linewidth}P{0.35\linewidth}@{}}
\toprule\textbf{Source / inspected window} & \textbf{Useful target} &
\textbf{Required qualification or remaining work}\\\midrule
\endhead
Anderson, PDF pp.~150--154 &
Oblique-shock conservation, tangential continuity, weak/strong roots and
attached limits. &
Verification of a calorically perfect operator; not an internal
shock-location or reacting-nozzle experiment.\\
Anderson, PDF pp.~665--668 &
General-EOS shock closure and upstream-state dependence. &
Air examples; match gas mode and reference definitions. An EOS compared
with itself is not an independent validation.\\
NASA/JHU source-port and prior Nozzle-B records &
Starting data, topology, constant-gas compatibility and resolved
inviscid numerical profiles. &
Retain differing source/model boundaries. Existing software acceptance
does not validate added gas/BL/shock models.\\
JPL TR 32-415, contents and Appendix B route &
Nozzle pressure and heat-transfer tables, calibration/instrumentation
discussion. &
Original table values, geometry and reduction/uncertainty still need a
bounded independently checked extraction \cite{back}.\\
NASA TP-2725, PDF pp.~3--5 &
1030:1 GO2/GH2 nozzle thrust, pressure and drag decrement. &
Campaign-specific inputs and uncertainty; prediction discrepancies in
the paper are not new repository results \cite{tp2725}.\\
NASA TP-2726, PDF pp.~3--5 &
Companion heat-transfer evidence at \(P_c=2413\) kPa and
O/F \(2.78\)--\(5.49\). &
Laminar/turbulent predictions differ materially; acceleration can
suppress turbulence. Do not assume a universal turbulent default
\cite{tp2726}.\\
NASA TP-3576, PDF pp.~3--4 &
High-area-ratio/high-pressure performance campaign. &
Reconstruct exact geometry, propellant and conditions before adding
a benchmark \cite{tp3576}.\\
\"Ostlund, PDF pp.~51--52, 70, 78--82 &
Onset scatter, FSS/RSS history and side-load definitions. &
Case-specific transitions; require matching wall/history and
appropriate pressure/load observables \cite{ostlund}.\\
Wang, PDF pp.~2, 5, 7 &
Reacting/cooled-wall model and heat/performance comparisons. &
Specific H2/O2 mechanism and cases; compare independent data rather
than merely matching the same numerical model \cite{wang}.\\
RBCC heat-transfer report, PDF p.~3 &
Wall-temperature/inverse heat-transfer methodology. &
It uses Bartz in a reduction step. Inferred heat flux is not automatically
an independent Bartz-validation target \cite{rbcc}.\\
\bottomrule
\end{longtable}

For attached BL work, the TP-2726 finding that laminar estimates can
agree much better than turbulent ones is a warning about model selection,
not a new universal numerical limit. The companion thrust and heat
reports share a campaign and must not be counted as independent repeated
experiments.

For separation/side loads, obtain wall-pressure histories, separation
definitions, circumferential measurements or direct loads and the
measurement/structural transfer model. An axisymmetric pressure curve
or an onset flag is insufficient for asymmetric transient loads.

\subsection{Numerical verification, model validation and optimization evidence}
The plan needs three complementary evidence classes:
\begin{itemize}
\item \textbf{Verification:} equations/discretization, conservative
jumps, property identities, manufactured or analytic limits, mesh and
gradient convergence.
\item \textbf{Model validation:} independent source-matched physical
observables with uncertainty and applicability; calibration data kept
separate from validation holdouts.
\item \textbf{Optimization evidence:} correct total sensitivities,
actual constraint enforcement, same-model state/geometry, converged
reoptimization/refinement and documented objective scope.
\end{itemize}
Passing any one class does not establish the other two. Improved nozzle
analysis also does not validate combustion, pumps, materials, manufacturing
or hardware.

\section{Open decisions and a concrete next work package}
\label{sec:next}
\subsection{Decisions to settle before implementing broader physics}
The initial target domain should name propellants/composition model,
pressure/altitude range, nozzle families, area/length/curvature ranges,
wall thermal conditions and Reynolds/transition assumptions. Decide
whether the first default is attached calorically perfect or requires
frozen variable-property flow. Specify acceptable runtime for interactive
and batch MDO, independent thermodynamic/flow references, and a
separated/mixed-type solver route where needed.

Source gaps remain explicit: original KL coefficients/start-line
validation; complete nozzle-specific attached-BL closures; corroborated
reacting rotational/axisymmetric source terms; selected reaction/transport
data; general-gas fitted-shock topology/pressure cases; and
geometry-complete separated-flow/load datasets.

\subsection{Recommended first implementation package}
The first bounded package should deliver C0 plus the N0/N1 design:
a versioned flow/performance contract, a prescribed-wall perfect-gas
primal topology reference, and an explicit total-derivative specification.
In parallel, T0 can expose the existing frozen-property model without
changing its physical meaning. The subsequent implementation should
demonstrate M0 on a restricted accepted case before changing any default.

This package addresses the current integration bottleneck: making the
engine optimize the same accepted nozzle flow whose geometry and
performance are reported. It produces a reviewable foundation for
variable gas, viscous interaction and fitted shocks, while retaining
explicit limits on unsupported regimes.

\subsection{Evidence package and status of this study}
The companion JSON files preserve code symbols/hashes, local literature
identities and inspected page windows, proposed milestones, and unresolved
source discrepancies. Local research reports are archived as supporting
evidence, with no bulk changes to the knowledge graph. The primary
differentiation paper and software documentation are explicitly external
supplements; no local document ID or hash was invented for them.

This study changed documentation only. Source-code presence and named
tests were inspected, and the previous accepted profile's source hashes
were checked. No new multi-physics implementation, fresh solver run,
default switch or end-to-end MDO improvement is claimed.
